\documentclass[14pt]{extarticle}
% ---------- PACKAGES ----------
\usepackage{amsmath, amssymb}
\usepackage{geometry}
\usepackage{setspace}
\usepackage{titlesec}
\usepackage{xcolor}
\usepackage{helvet}
\usepackage{enumitem}
\usepackage{lmodern}
\makeatletter
\IfFileExists{../common-things/reflectionpage.tex}{%
  \input{../common-things/reflectionpage.tex}%
}{%
  \IfFileExists{common-things/reflectionpage.tex}{%
    \input{common-things/reflectionpage.tex}%
  }{%
    \PackageError{\jobname}{Missing reflectionpage.tex file}{%
      Place reflectionpage.tex inside a common-things/ directory next to this unit\@.%
    }%
  }%
}
\makeatother
% ---------- PAGE SETUP ----------
\geometry{margin=1in}
\setstretch{1.5}
\setlength{\parindent}{0pt}
\setlength{\parskip}{0.75em}
\setlist{itemsep=0.75\baselineskip, topsep=0.5\baselineskip}
\titleformat{\section}{\normalfont\Large\bfseries}{\thesection}{1em}{}
\titleformat{\subsection}{\normalfont\large\bfseries}{\thesubsection}{1em}{}
\renewcommand{\familydefault}{\sfdefault}
\renewcommand{\emph}[1]{\textbf{#1}}

\begin{document}
\raggedright
\pagenumbering{gobble}

\begin{center}
    \LARGE \textbf{Unit 8: Geometry and Trigonometry} \\[6pt]
    \Large \textbf{Topic 4: The Pythagorean Theorem and Special Right Triangles}
\end{center}

\vspace{1em}

\section*{Concept Summary}

The \textbf{Pythagorean Theorem} states that in any right triangle with legs \(a\) and \(b\) and hypotenuse \(c\):
\[
a^2 + b^2 = c^2
\]
The hypotenuse is always the longest side and is always opposite the right angle. This theorem lets you find any one side of a right triangle when the other two are known. It also works in reverse: if three sides satisfy \(a^2 + b^2 = c^2\), the triangle is a right triangle.

Certain sets of whole numbers satisfy the Pythagorean Theorem and appear frequently on the SAT. These are called \textbf{Pythagorean triples}. The most common are:
\[
3\text{-}4\text{-}5, \quad 5\text{-}12\text{-}13, \quad 8\text{-}15\text{-}17, \quad 7\text{-}24\text{-}25
\]
Any constant multiple of a Pythagorean triple is also a triple. For example, \(3\text{-}4\text{-}5\) gives \(6\text{-}8\text{-}10\), \(9\text{-}12\text{-}15\), \(15\text{-}20\text{-}25\), and so on. Recognizing these triples saves time because you can identify the missing side by inspection instead of computing squares and square roots.

Two families of right triangles have side ratios worth memorizing because the SAT provides them on the reference sheet and tests them often:

A \textbf{45-45-90 triangle} (isosceles right triangle) has sides in the ratio \(1 : 1 : \sqrt{2}\). If each leg has length \(s\), the hypotenuse is \(s\sqrt{2}\). Conversely, if the hypotenuse is \(h\), each leg is \(\dfrac{h}{\sqrt{2}} = \dfrac{h\sqrt{2}}{2}\).

A \textbf{30-60-90 triangle} has sides in the ratio \(1 : \sqrt{3} : 2\). The shortest side (opposite the \(30^\circ\) angle) has length \(s\), the side opposite the \(60^\circ\) angle has length \(s\sqrt{3}\), and the hypotenuse (opposite the \(90^\circ\) angle) has length \(2s\). The most common mistake is mixing up which side gets the \(\sqrt{3}\) factor. Remember: the \(\sqrt{3}\) side is the middle-length side, between the shortest side and the hypotenuse.

\section*{Core Skills}
\begin{itemize}
  \item Use the Pythagorean Theorem to find a missing side of a right triangle.
  \item Recognize common Pythagorean triples and their multiples to save computation time.
  \item Apply the side ratios of 45-45-90 and 30-60-90 triangles to find missing sides.
  \item Use the converse of the Pythagorean Theorem to determine whether a triangle is a right triangle.
  \item Solve multi-step problems that combine the Pythagorean Theorem with area, perimeter, or other geometry concepts.
\end{itemize}

\section*{Example 1: Basic Pythagorean Theorem}

A right triangle has legs of length 9 and 12. What is the length of the hypotenuse?

\textbf{Step 1:} Apply the Pythagorean Theorem.
\[
c^2 = 9^2 + 12^2 = 81 + 144 = 225
\]

\textbf{Step 2:} Take the square root.
\[
c = \sqrt{225} = 15
\]

(You could also recognize this as the \(3\text{-}4\text{-}5\) triple scaled by 3: \(9\text{-}12\text{-}15\).)

The hypotenuse is \(\boxed{15}\).

\section*{Example 2: Finding a Leg}

A right triangle has a hypotenuse of 26 and one leg of 10. What is the length of the other leg?

\textbf{Step 1:} Let the unknown leg be \(a\).
\[
a^2 + 10^2 = 26^2 \implies a^2 + 100 = 676
\]

\textbf{Step 2:} Solve for \(a\).
\[
a^2 = 576 \implies a = 24
\]

(This is the \(5\text{-}12\text{-}13\) triple scaled by 2: \(10\text{-}24\text{-}26\).)

The other leg is \(\boxed{24}\).

\section*{Example 3: 45-45-90 Triangle}

In a 45-45-90 triangle, the hypotenuse has length 10. What is the length of each leg?

\textbf{Step 1:} In a 45-45-90 triangle, the hypotenuse is \(s\sqrt{2}\), where \(s\) is the leg length.
\[
s\sqrt{2} = 10 \implies s = \frac{10}{\sqrt{2}} = \frac{10\sqrt{2}}{2} = 5\sqrt{2}
\]

Each leg has length \(\boxed{5\sqrt{2}}\).

\section*{Example 4: 30-60-90 Triangle}

In a 30-60-90 triangle, the side opposite the \(30^\circ\) angle has length 7. What are the lengths of the other two sides?

\textbf{Step 1:} The shortest side (opposite \(30^\circ\)) is \(s = 7\).

\textbf{Step 2:} The side opposite \(60^\circ\) is \(s\sqrt{3} = 7\sqrt{3}\).

\textbf{Step 3:} The hypotenuse is \(2s = 14\).

The side opposite \(60^\circ\) is \(\boxed{7\sqrt{3}}\) and the hypotenuse is \(\boxed{14}\).

\section*{Example 5: Converse of the Pythagorean Theorem}

A triangle has sides of length 11, 60, and 61. Is it a right triangle?

\textbf{Step 1:} Check whether the square of the longest side equals the sum of the squares of the other two.
\[
11^2 + 60^2 = 121 + 3600 = 3721
\]
\[
61^2 = 3721
\]

Since \(11^2 + 60^2 = 61^2\), the triangle is a right triangle.

\(\boxed{\text{Yes}}\)

\section*{Example 6: Multi-Step Application}

A rectangular room measures 8 feet by 15 feet. What is the length, in feet, of a diagonal of the room?

\textbf{Step 1:} A diagonal of a rectangle divides it into two right triangles. The legs are the length and width of the room.

\textbf{Step 2:} Apply the Pythagorean Theorem.
\[
d^2 = 8^2 + 15^2 = 64 + 225 = 289 \implies d = 17
\]

(This is the \(8\text{-}15\text{-}17\) Pythagorean triple.)

The diagonal is \(\boxed{17}\) feet.

\section*{Key Takeaways}
\begin{itemize}
  \item Memorize the common Pythagorean triples (\(3\text{-}4\text{-}5\), \(5\text{-}12\text{-}13\), \(8\text{-}15\text{-}17\)) and watch for their multiples. Recognizing a triple instantly gives you the answer without any calculation.
  \item In a 45-45-90 triangle, multiply a leg by \(\sqrt{2}\) to get the hypotenuse, or divide the hypotenuse by \(\sqrt{2}\) to get a leg. In a 30-60-90 triangle, the sides are \(s\), \(s\sqrt{3}\), and \(2s\), with the \(\sqrt{3}\) always on the middle side (opposite \(60^\circ\)).
  \item The most common error with 30-60-90 triangles is applying the \(\sqrt{3}\) to the wrong side. The hypotenuse is always exactly twice the shortest side, never involving \(\sqrt{3}\).
  \item When a problem gives you a rectangle, square, or equilateral triangle and asks for a diagonal or height, you are almost certainly using the Pythagorean Theorem or a special right triangle.
\end{itemize}

\newpage

\section*{Practice Questions: The Pythagorean Theorem and Special Right Triangles}

\subsection*{Part A: Pythagorean Theorem Basics}
\begin{enumerate}
  \item A right triangle has legs of length 5 and 12. What is the length of the hypotenuse?
  \item A right triangle has a hypotenuse of 25 and one leg of 7. What is the length of the other leg?
  \item A right triangle has legs of length 6 and 8. What is the length of the hypotenuse?
  \item A right triangle has a hypotenuse of 20 and one leg of 16. What is the length of the other leg?
  \item A right triangle has legs of length 15 and 20. What is the length of the hypotenuse?
\end{enumerate}

\subsection*{Part B: Special Right Triangles}
\begin{enumerate}
  \item In a 45-45-90 triangle, each leg has length 6. What is the length of the hypotenuse?
  \item In a 45-45-90 triangle, the hypotenuse has length 14. What is the length of each leg? Express your answer in simplified radical form.
  \item In a 30-60-90 triangle, the hypotenuse has length 18. What is the length of the side opposite the \(30^\circ\) angle?
  \item In a 30-60-90 triangle, the side opposite the \(60^\circ\) angle has length \(5\sqrt{3}\). What is the length of the hypotenuse?
  \item An equilateral triangle has side length 12. What is the height of the triangle? Express your answer in simplified radical form.
\end{enumerate}

\subsection*{Part C: Converse and Classification}
\begin{enumerate}
  \item A triangle has sides of length 9, 40, and 41. Is it a right triangle?
  \item A triangle has sides of length 7, 10, and 13. Is it a right triangle?
  \item A triangle has sides of length 20, 21, and 29. Is it a right triangle?
  \item The lengths 14, 48, and 50 form a triangle. Is this triangle a right triangle?
  \item A triangle has sides of length \(2k\), \(k^2 - 1\), and \(k^2 + 1\), where \(k > 1\). Show that these always form a right triangle by verifying the Pythagorean relationship. For \(k = 4\), what are the three side lengths?
\end{enumerate}

\subsection*{Part D: Multi-Step Problems}
\begin{enumerate}
  \item A right triangle has legs of length 9 and 40. What is the area of the triangle?
  \item A square has a diagonal of length \(10\sqrt{2}\). What is the side length of the square?
  \item A 12-foot ladder leans against a wall with its base 5 feet from the wall. How high up the wall does the ladder reach, in feet?
  \item The two diagonals of a rhombus have lengths 10 and 24. What is the length of each side of the rhombus?
  \item A right triangle has a hypotenuse of 50 and one leg of 30. What is the perimeter of the triangle?
\end{enumerate}

\subsection*{Part E: SAT-Style Applications}
\begin{enumerate}
  \item A cable runs from the top of a 45-foot utility pole to a point on the ground 28 feet from the base of the pole. What is the length of the cable, in feet?
  \item A rectangular television screen has a width of 48 inches and a diagonal of 52 inches. What is the height of the screen, in inches?
  \item A ramp in the shape of a 30-60-90 triangle rises to a platform that is 4 feet above the ground. What is the length, in feet, of the ramp (the hypotenuse)?
  \item A square courtyard has an area of 200 square meters. A path runs diagonally from one corner to the opposite corner. What is the length of the path, in meters?
  \item A ship sails 24 miles due north and then 10 miles due east. What is the shortest distance, in miles, between the ship's starting point and its final position?
\end{enumerate}

\newpage

\section*{Answer Key and Solutions: The Pythagorean Theorem and Special Right Triangles}

\subsection*{Part A Solutions: Pythagorean Theorem Basics}
\begin{enumerate}
  \item \(c^2 = 5^2 + 12^2 = 25 + 144 = 169\), so \(c = 13\). (This is the \(5\text{-}12\text{-}13\) triple.)

  \(\boxed{13}\)

  \item \(a^2 + 7^2 = 25^2 \implies a^2 = 625 - 49 = 576 \implies a = 24\). (This is the \(7\text{-}24\text{-}25\) triple.)

  \(\boxed{24}\)

  \item \(c^2 = 6^2 + 8^2 = 36 + 64 = 100\), so \(c = 10\). (This is the \(3\text{-}4\text{-}5\) triple scaled by 2.)

  \(\boxed{10}\)

  \item \(a^2 + 16^2 = 20^2 \implies a^2 = 400 - 256 = 144 \implies a = 12\). (This is the \(3\text{-}4\text{-}5\) triple scaled by 4: \(12\text{-}16\text{-}20\).)

  \(\boxed{12}\)

  \item \(c^2 = 15^2 + 20^2 = 225 + 400 = 625\), so \(c = 25\). (This is the \(3\text{-}4\text{-}5\) triple scaled by 5.)

  \(\boxed{25}\)
\end{enumerate}

\subsection*{Part B Solutions: Special Right Triangles}
\begin{enumerate}
  \item In a 45-45-90 triangle with leg \(s = 6\), the hypotenuse is \(s\sqrt{2} = 6\sqrt{2}\).

  \(\boxed{6\sqrt{2}}\)

  \item The hypotenuse is \(s\sqrt{2} = 14\), so \(s = \dfrac{14}{\sqrt{2}} = \dfrac{14\sqrt{2}}{2} = 7\sqrt{2}\).

  \(\boxed{7\sqrt{2}}\)

  \item In a 30-60-90 triangle, the hypotenuse is \(2s\), where \(s\) is the shortest side (opposite \(30^\circ\)).
  \[
  2s = 18 \implies s = 9
  \]

  \(\boxed{9}\)

  \item The side opposite \(60^\circ\) is \(s\sqrt{3}\), so \(s\sqrt{3} = 5\sqrt{3}\), which gives \(s = 5\). The hypotenuse is \(2s = 10\).

  \(\boxed{10}\)

  \item The height of an equilateral triangle splits it into two 30-60-90 triangles. The base of each is \(\dfrac{12}{2} = 6\) (the side opposite \(30^\circ\)), and the height is the side opposite \(60^\circ\):
  \[
  h = 6\sqrt{3}
  \]

  \(\boxed{6\sqrt{3}}\)
\end{enumerate}

\subsection*{Part C Solutions: Converse and Classification}
\begin{enumerate}
  \item Check: \(9^2 + 40^2 = 81 + 1600 = 1681\) and \(41^2 = 1681\). Since \(9^2 + 40^2 = 41^2\), it is a right triangle.

  \(\boxed{\text{Yes}}\)

  \item Check: \(7^2 + 10^2 = 49 + 100 = 149\) and \(13^2 = 169\). Since \(149 \neq 169\), it is not a right triangle.

  \(\boxed{\text{No}}\)

  \item Check: \(20^2 + 21^2 = 400 + 441 = 841\) and \(29^2 = 841\). Since \(20^2 + 21^2 = 29^2\), it is a right triangle.

  \(\boxed{\text{Yes}}\)

  \item Check: \(14^2 + 48^2 = 196 + 2304 = 2500\) and \(50^2 = 2500\). Since \(14^2 + 48^2 = 50^2\), it is a right triangle. (This is the \(7\text{-}24\text{-}25\) triple scaled by 2.)

  \(\boxed{\text{Yes}}\)

  \item We need to verify that \((2k)^2 + (k^2 - 1)^2 = (k^2 + 1)^2\).
  \[
  (2k)^2 + (k^2 - 1)^2 = 4k^2 + k^4 - 2k^2 + 1 = k^4 + 2k^2 + 1 = (k^2 + 1)^2
  \]
  The identity holds for all \(k > 1\), so these always form a right triangle. For \(k = 4\): the sides are \(2(4) = 8\), \(16 - 1 = 15\), and \(16 + 1 = 17\).

  \(\boxed{8, 15, 17}\)
\end{enumerate}

\subsection*{Part D Solutions: Multi-Step Problems}
\begin{enumerate}
  \item The area of a right triangle is \(\dfrac{1}{2} \cdot \text{leg}_1 \cdot \text{leg}_2 = \dfrac{1}{2}(9)(40) = 180\).

  \(\boxed{180}\)

  \item A square's diagonal and side are related by the 45-45-90 ratio: diagonal \(= s\sqrt{2}\). So \(s\sqrt{2} = 10\sqrt{2}\), which gives \(s = 10\).

  \(\boxed{10}\)

  \item The ladder, wall, and ground form a right triangle with hypotenuse 12 and one leg 5.
  \[
  h^2 + 5^2 = 12^2 \implies h^2 = 144 - 25 = 119 \implies h = \sqrt{119}
  \]
  Since \(\sqrt{119}\) is not a clean number, let us verify: this is not a standard Pythagorean triple. The ladder reaches \(\sqrt{119}\) feet up the wall. As a decimal, this is approximately \(10.9\) feet.

  \(\boxed{\sqrt{119}}\)

  \item The diagonals of a rhombus bisect each other at right angles. Each half-diagonal has length \(\dfrac{10}{2} = 5\) and \(\dfrac{24}{2} = 12\). Each side of the rhombus is the hypotenuse of a right triangle with legs 5 and 12.
  \[
  s^2 = 5^2 + 12^2 = 25 + 144 = 169 \implies s = 13
  \]

  \(\boxed{13}\)

  \item First find the missing leg: \(a^2 + 30^2 = 50^2 \implies a^2 = 2500 - 900 = 1600 \implies a = 40\). (This is the \(3\text{-}4\text{-}5\) triple scaled by 10.)
  \[
  P = 30 + 40 + 50 = 120
  \]

  \(\boxed{120}\)
\end{enumerate}

\subsection*{Part E Solutions: SAT-Style Applications}
\begin{enumerate}
  \item The pole, ground, and cable form a right triangle with legs 45 and 28.
  \[
  c^2 = 45^2 + 28^2 = 2025 + 784 = 2809 \implies c = 53
  \]

  The cable is \(\boxed{53}\) feet long.

  \item Let the height be \(h\). Then \(48^2 + h^2 = 52^2\).
  \[
  h^2 = 2704 - 2304 = 400 \implies h = 20
  \]

  The height is \(\boxed{20}\) inches.

  \item The platform height of 4 feet is the side opposite the \(30^\circ\) angle (the shortest side) in the 30-60-90 triangle. The hypotenuse is twice the shortest side.
  \[
  \text{ramp length} = 2 \times 4 = 8
  \]

  The ramp is \(\boxed{8}\) feet long.

  \item The side length of the square is \(s = \sqrt{200} = 10\sqrt{2}\). The diagonal of a square is \(s\sqrt{2} = 10\sqrt{2} \cdot \sqrt{2} = 10 \cdot 2 = 20\).

  \(\boxed{20}\)

  \item The north and east directions are perpendicular, forming a right triangle with legs 24 and 10.
  \[
  d^2 = 24^2 + 10^2 = 576 + 100 = 676 \implies d = 26
  \]

  The shortest distance is \(\boxed{26}\) miles.
\end{enumerate}

\ReflectionPage{Unit 8, Topic 4}{https://www.scotchegg.co}

\end{document}
